\documentclass[11pt,letterpaper]{article}
\usepackage{fontspec}
\usepackage[spanish,es-nodecimaldot,es-noindentfirst,es-noshorthands]{babel}
\usepackage[margin=2.1cm]{geometry}
\usepackage{amsmath,amssymb}
\usepackage{longtable,booktabs,array}
\usepackage{xcolor}
\usepackage{enumitem}
\usepackage{titlesec}
\usepackage[hidelinks]{hyperref}
\usepackage{tcolorbox}

\definecolor{agua}{RGB}{8,81,156}
\definecolor{celeste}{RGB}{224,245,252}
\definecolor{gris}{RGB}{90,90,90}
\titleformat{\section}{\Large\bfseries\color{agua}}{\thesection}{0.6em}{}
\titleformat{\subsection}{\large\bfseries\color{agua}}{\thesubsection}{0.6em}{}
\setlist{itemsep=2pt,topsep=4pt}
\setlength{\parskip}{4pt}
\setlength{\emergencystretch}{3em}

\newtcolorbox{decir}{colback=celeste,colframe=agua,boxrule=0.5pt,arc=2pt,left=6pt,right=6pt,top=4pt,bottom=4pt}
\newcommand{\mostrar}[1]{\par\noindent{\small\color{gris}\textbf{En pantalla:} #1}\par}

\begin{document}

\begin{center}
{\LARGE\bfseries\color{agua} Simulador de inundación por lluvia}\\[4pt]
{\large TerraPulse AI — guion de presentación, parámetros y ecuaciones}\\[4pt]
{\small Rama \texttt{feature/flood-ui-dg2} · 21 de noviembre de 2026}
\end{center}

\section{Qué es y para qué sirve}

El simulador responde una pregunta concreta: \textbf{si cae cierta lluvia sobre una zona, ¿hacia dónde se mueve el agua y dónde se acumula?}
No pronostica el clima ni reproduce un evento observado: simula \emph{escenarios} de lluvia (inventada o de diseño) para prever qué zonas se mojan, por dónde escurre el agua y dónde se encharca.

\begin{itemize}
  \item \textbf{Propósito:} explorar escenarios de lluvia y ver el recorrido del agua sobre el terreno real, en cualquier parte del mundo y solo con datos abiertos.
  \item \textbf{Motor:} LISFLOOD-FP 8.0.3, un solver hidráulico de código abierto que resuelve las ecuaciones de aguas someras en 2D. Se usan dos métodos: DG2 (Galerkin discontinuo de 2.º orden, más preciso) y FV1 (volúmenes finitos de 1.er orden, más estable en terreno empinado). Si el método elegido falla, el servidor prueba el otro automáticamente y lo avisa.
  \item \textbf{Datos de entrada (todos abiertos y globales):} terreno Copernicus GLO-30 o NASADEM (NASA Earthdata/ASF), cobertura ESA WorldCover, suelos HYSOGs250m (Earthdata), humedad NISAR SME2 o SMAP L4, lluvia de diseño PXR-4, mareas EOT20, fondo marino GEBCO y lagos HydroLAKES.
  \item \textbf{Resultado:} una animación cuadro por cuadro (mientras corre y al final), el mapa de profundidad máxima, la distribución por clases de profundidad, el área inundada y un balance de agua que dice cuánta lluvia cayó, cuánta se infiltró, cuánta quedó en láminas, cuánta se encharcó y cuánta salió del dominio.
  \item \textbf{Lo que no es:} los resultados son \emph{exploratorios}. Todo se marca \texttt{synthetic} y \texttt{hydraulic\_accuracy\_\allowbreak validated:\ false}: no hay validación hidráulica contra un evento real.
\end{itemize}

\section{Cómo funciona, en seis pasos}

\begin{enumerate}
  \item \textbf{Zona.} El usuario dibuja o elige una ROI. Se usa como dominio de cálculo (no es automáticamente una cuenca). El presupuesto actual admite hasta 3000 celdas.
  \item \textbf{Terreno.} Se descarga el modelo de elevación de 30 m, se reproyecta a una malla métrica (UTM) y todas las alturas se llevan a un solo datum vertical, EGM2008, el mismo del nivel del mar.
  \item \textbf{Superficie.} Cada celda recibe su rugosidad de Manning según la cobertura (bosque, pasto, ciudad, agua…) y su número de curva (cobertura + tipo de suelo + humedad previa), que decide qué parte de la lluvia se infiltra.
  \item \textbf{Lluvia.} Sintética (total y duración que define el usuario, con un perfil aleatorio reproducible) o tormenta de diseño (la lluvia de un periodo de retorno, por ejemplo 100 años, en ese punto del mundo). La lluvia neta de cada celda entra al solver como un cubo NetCDF.
  \item \textbf{Solver.} LISFLOOD-FP resuelve el flujo con DG2 o FV1. El servidor vigila las profundidades negativas por volumen y, si el método falla, cambia al otro.
  \item \textbf{Resultados.} Se convierten las salidas a mapas (COG y PNG) y a la animación; se calcula el balance de agua y se muestran los avisos (zona urbana, artefactos numéricos, cambio de método).
\end{enumerate}

\section{Guion de presentación (8–10 minutos)}

\subsection*{1. Apertura — el problema (1 min)}
\begin{decir}
``Cuando llueve fuerte, la pregunta que importa no es cuánto llovió, sino a dónde se va el agua. Las calles, las laderas y los cañones deciden dónde se encharca. Nuestro simulador toma el terreno real de cualquier lugar del mundo, le pone encima una lluvia y calcula, celda por celda, cómo escurre y dónde se acumula.''
\end{decir}
\mostrar{el mapa con la zona del Cañón del Sumidero o de Tuxtla Gutiérrez.}

\subsection*{2. Qué hace y qué no hace (1 min)}
\begin{decir}
``Es una herramienta de escenarios: ¿qué pasaría si cae la lluvia de una tormenta de 100 años durante 3 horas? No es un pronóstico ni está validado contra un evento real; por eso todos los resultados dicen `no validado'. Lo que sí garantiza es que el agua se conserva: lo que entra como lluvia, se infiltra, queda en el terreno o sale del dominio.''
\end{decir}

\subsection*{3. Los datos — todo abierto y global (1.5 min)}
\begin{decir}
``Usamos solo datos abiertos para que funcione en cualquier parte del mundo. El terreno viene de Copernicus GLO-30 a través de ASF y NASA Earthdata. La rugosidad del suelo sale de ESA WorldCover; los tipos de suelo, de HYSOGs de Earthdata; la humedad previa, de NISAR SME2 o de SMAP. La lluvia de diseño sale de PXR-4, unas curvas de intensidad-duración-frecuencia globales. Todas las alturas están en un mismo datum, EGM2008, para que terreno y nivel del mar sean comparables.''
\end{decir}

\subsection*{4. Demostración en vivo — configurar (2 min)}
\begin{decir}
``Elijo la zona en el mapa. Con `Cargar valores de ejemplo' se llena el panel: una tormenta de diseño de 100 años y 3 horas, cuatro horas de simulación y un cuadro cada 5 minutos. El terreno es Copernicus a 30 metros, la rugosidad sale de la cobertura y la infiltración del número de curva con la humedad del suelo automática. Todas las fronteras están libres, así que el agua puede salir. El método es DG2; si falla, el servidor pasa solo a FV1 y nos avisa.''
\end{decir}
\mostrar{el panel paso a paso: Zona → Lluvia → Terreno y fronteras → Opciones avanzadas → Confirmar y simular.}

\subsection*{5. Mientras calcula — la animación como barra de progreso (1.5 min)}
\begin{decir}
``Mientras el solver trabaja vemos el progreso: el tiempo simulado avanza y aparecen los cuadros que ya calculó. Las gotas y el velo muestran que está lloviendo sobre todo el dominio, y el indicador dice la intensidad y cuánto lleva acumulado. El celeste claro es la superficie mojada: desde medio milímetro ya se ve que llovió en toda la zona. Conforme avanza, el agua se junta en las partes bajas y el azul se vuelve más fuerte.''
\end{decir}

\subsection*{6. Resultados — dónde se fue el agua (1.5 min)}
\begin{decir}
``Al final tenemos la reproducción completa, el mapa de profundidad máxima y el balance de agua. Por ejemplo, en el Sumidero cayeron unos 9\,000 m³ de lluvia; cerca del 70\% se infiltró, y de lo que escurrió casi todo terminó encharcado en dos fondos de valle con hasta un metro de profundidad. Las láminas delgadas en las laderas son de milímetros: el agua pasa por ahí, pero se acumula abajo.''
\end{decir}
\mostrar{la gráfica de distribución por clases y el balance de agua.}

\subsection*{7. Límites y cierre (1 min)}
\begin{decir}
``Somos honestos con los límites: no modelamos ríos que entran desde fuera, ni drenaje urbano, ni evaporación; el terreno de 30 metros incluye árboles y edificios, y las curvas de lluvia globales suavizan las tormentas tropicales cortas. Por eso es una herramienta para explorar escenarios, no para decidir una evacuación. El siguiente paso es validarlo contra una inundación observada por satélite.''
\end{decir}

\section{Parámetros del panel}

La tabla sigue el orden del panel. ``Por defecto/ejemplo'' indica el valor que pone ``Cargar valores de ejemplo''; no son recomendaciones físicas.

{\small
\renewcommand{\arraystretch}{1.25}
\begin{longtable}{>{\raggedright\arraybackslash}p{2.1cm}>{\raggedright\arraybackslash}p{3.0cm}>{\raggedright\arraybackslash}p{5.3cm}>{\raggedright\arraybackslash}p{5.2cm}}
\toprule
\textbf{Sección} & \textbf{Parámetro} & \textbf{Qué hace} & \textbf{Por qué se necesita}\\
\midrule
\endfirsthead
\toprule
\textbf{Sección} & \textbf{Parámetro} & \textbf{Qué hace} & \textbf{Por qué se necesita}\\
\midrule
\endhead
\bottomrule
\endfoot
% ---------------- Zona
1. Zona & ROI (polígono) & Define el área donde se calcula el flujo. & Sin un dominio no hay malla; el tamaño de la zona y la resolución deciden cuántas celdas hay (máx. 3000).\\
 & Declaración ``uso la ROI como dominio'' & El usuario acepta que la ROI es un dominio de cálculo, no una cuenca. & Una ROI rectangular corta cuencas: el agua que vendría de fuera no entra. Hay que saberlo para interpretar.\\
 & Otro dominio (GeoJSON) & Permite un dominio distinto que cubra la ROI. & Sirve para ampliar el área o centrarla en un valle.\\
% ---------------- Lluvia
2. Lluvia & Tipo de lluvia & Sintética (la define el usuario) o tormenta de diseño (PXR-4). & Elegir entre probar un escenario libre o una lluvia con probabilidad conocida en ese lugar.\\
 & Periodo de retorno (años) & 2, 5, 10, 25, 50 o 100 años. Probabilidad anual $1/T$ de superar esa lluvia. Ejemplo: 100. & Traduce ``qué tan extrema'' es la tormenta a una cantidad de lluvia real del sitio.\\
 & Duración de la tormenta (h) & Horas de lluvia, en múltiplos de 10 min (máx. 24 h). Ejemplo: 3 h. & La misma probabilidad da más milímetros si dura más; la intensidad pico depende de la duración.\\
 & Lluvia hipotética (mm) & Total que cae sobre la zona (modo sintético). & Es el volumen de agua del escenario.\\
 & Duración de la lluvia (s) & Cuánto tiempo llueve (modo sintético). & Con el total define la intensidad media (mm/h).\\
 & Número de intervalos & Divide la lluvia en tramos de igual duración con intensidades distintas; el total se conserva. & Una lluvia real no es constante; los tramos dan picos y pausas.\\
 & Semilla & Entero que fija el perfil aleatorio. & Reproducibilidad: la misma semilla repite exactamente la misma lluvia.\\
% ---------------- Terreno y fronteras
3. Terreno y fronteras & Terreno & Copernicus GLO-30 (recomendado, global) o NASADEM (56°\,S–60°\,N), ambos de 30 m vía Earthdata. & La forma del terreno decide hacia dónde corre el agua; es la entrada más importante.\\
 & Escenario rápido & ``Tierra adentro'' (todo libre) o ``Mar al norte/este/sur/oeste''. & Atajo para configurar las fronteras sin errores.\\
 & Frontera por lado (N, E, S, O) & Cerrada (pared), libre (el agua sale), nivel fijo (mar) o marea EOT20. & Define qué pasa con el agua en el borde del dominio; un borde mal puesto acumula o pierde agua artificialmente.\\
 & Inicio del evento (UTC) & Fecha y hora de inicio (solo con marea). & La marea depende del instante; se calcula de inicio a inicio + duración.\\
 & Nivel del mar fijo (m, EGM2008) & Altura del mar en los lados de nivel fijo. & Escenarios costeros: el mar alto impide que el agua salga.\\
 & Iniciar con el mar lleno & Llena al inicio las celdas bajo el nivel del mar conectadas al lado de mar. & En DG2 un lado de nivel fijo no llena un dominio seco por sí solo.\\
 & Fondo marino GEBCO & Usa la batimetría GEBCO\_2026 en las celdas de mar. & El terreno de 30 m no tiene fondo marino; sin él la profundidad inicial del mar sería cero.\\
 & Presión atmosférica (hPa) & Sube el nivel del mar $\approx$1 cm por hPa bajo 1013.25. & Representa el efecto de barómetro inverso de una tormenta (opcional).\\
 & Rugosidad de Manning & Uniforme (un valor $n$) o según cobertura (un $n$ por celda de WorldCover). & La rugosidad frena el agua: un bosque ($n\approx0.10$) la frena más que un pastizal ($n\approx0.035$).\\
 & Manning $n$ (s/m$^{1/3}$) & Valor único cuando el modo es uniforme. Ejemplo: 0.05. & Controla la velocidad del flujo en todo el dominio.\\
 & Infiltración & No modelada, coeficiente de escorrentía (0.65–0.92) o número de curva por celda. & Solo una parte de la lluvia escurre; el resto entra al suelo. Sin esto se sobreestima el agua.\\
 & Humedad previa del suelo & Automática (NISAR SME2 o SMAP L4), I seco, II medio o III húmedo. & Un suelo ya mojado infiltra menos y escurre más.\\
 & Lagos & No incluir, o HydroLAKES (lagos llenos al inicio). & Un lago vacío al inicio se tragaría la lluvia de toda la zona.\\
% ---------------- Avanzadas
4. Opciones avanzadas & Método & DG2 (más preciso) o FV1 (más estable en terreno empinado). Respaldo automático al otro si falla. & DG2 resuelve mejor el detalle, pero en laderas muy empinadas puede generar profundidades negativas.\\
 & CRS métrico (EPSG) & Proyección en metros (Chiapas: EPSG:32615, UTM 15N). Se calcula sola. & El solver trabaja en metros; los grados no sirven para distancias.\\
 & Resolución (m) & Tamaño de celda. Ejemplo: 30 m. & Define el detalle y el número de celdas (y con ello el tiempo de cálculo).\\
 & Remuestreo & \texttt{nearest} conserva las alturas; \texttt{bilinear} las suaviza. & Al cambiar de malla hay que interpolar el terreno; \texttt{nearest} no inventa alturas.\\
 & Duración de simulación (s) & Tiempo simulado (máx. 86\,400 s). Ejemplo: 14\,400 s = 4 h. & Debe cubrir toda la lluvia y dejar tiempo extra para ver a dónde escurre el agua.\\
 & Intervalo de guardado (s) & Cada cuánto se guarda un cuadro. Ejemplo: 300 s (49 cuadros). & Es la resolución temporal de la animación (máx. 121 cuadros).\\
 & Paso inicial (s) & Primer paso de tiempo; luego el solver lo ajusta. Ejemplo: 1 s. & Arranque estable del esquema numérico.\\
 & Umbral de lámina delgada DG2 (m) & Elimina profundidades negativas en láminas finas. Por defecto desactivado. & Estabiliza DG2, pero cambia la solución; por eso queda registrado.\\
% ---------------- Confirmar
5. Confirmar y simular & Confirmación de hipótesis & El usuario confirma que entiende que el resultado es exploratorio. & Evita presentar un escenario como si fuera una predicción validada.\\
% ---------------- Resultados
Resultados & Reproducción (play, velocidad 0.5×–4×, repetir) & Anima los cuadros en orden; muestra el tiempo simulado. & Se ve cómo se van mojando y llenando las zonas.\\
 & Mostrar lluvia & Gotas, velo e indicador ``Lloviendo X mm/h · acumulado Y mm de Z mm''. & Hace visible que llueve en toda la zona y con qué intensidad.\\
 & Ver máximo & Mapa de la profundidad máxima en toda la corrida. & Resume las zonas que más se inundaron.\\
 & Clases de profundidad & 0.5–2 mm superficie mojada; 2–10 mm escurrimiento; 1–3 cm lámina delgada; desde 3 cm encharcada (hasta $\geq$2 m). & Separa ``pasó agua'' de ``se acumuló agua''.\\
 & Área inundada & Área con profundidad $\geq$ 0.03 m. & Indicador simple y comparable entre corridas.\\
 & Balance de agua & Lluvia bruta, infiltrada, neta, en láminas, encharcada y salida por los bordes (m³). & Demuestra que el agua se conserva y explica a dónde se fue.\\
 & Avisos & Zona urbana, artefactos numéricos, cambio de método. & Transparencia: el usuario sabe cuándo confiar menos en las profundidades.\\
\end{longtable}
}

\section{Ecuaciones}

\subsection{Flujo del agua: ecuaciones de aguas someras 2D}
Variables: profundidad $h$, velocidades $u,v$, terreno $z_b$, gravedad $g$, rugosidad $n$. LISFLOOD-FP DG2/FV1 resuelve la forma conservativa:
\begin{align}
\frac{\partial h}{\partial t}+\frac{\partial (hu)}{\partial x}+\frac{\partial (hv)}{\partial y} &= R_{\text{neta}} \label{eq:masa}\\
\frac{\partial (hu)}{\partial t}+\frac{\partial}{\partial x}\!\left(hu^2+\tfrac12 g h^2\right)+\frac{\partial (huv)}{\partial y} &= -g h\frac{\partial z_b}{\partial x}-g h S_{fx}\\
\frac{\partial (hv)}{\partial t}+\frac{\partial (huv)}{\partial x}+\frac{\partial}{\partial y}\!\left(hv^2+\tfrac12 g h^2\right) &= -g h\frac{\partial z_b}{\partial y}-g h S_{fy}
\end{align}
La ecuación \eqref{eq:masa} es la conservación de masa: el agua de una celda cambia por lo que entra y sale por sus lados más la lluvia neta $R_{\text{neta}}$ (m/s). Las otras dos son la conservación de momento: la pendiente del terreno acelera el agua y la fricción la frena.

\paragraph{Fricción de Manning.}
\begin{equation}
S_{fx}=\frac{n^2\,u\sqrt{u^2+v^2}}{h^{4/3}},\qquad S_{fy}=\frac{n^2\,v\sqrt{u^2+v^2}}{h^{4/3}}
\end{equation}
Con más rugosidad $n$ o menos profundidad $h$, la fricción es mayor: por eso las láminas delgadas se mueven lento.

\paragraph{Estabilidad (condición CFL).} El solver es explícito y ajusta su paso de tiempo para que el agua no cruce más de una celda por paso:
\begin{equation}
\Delta t \le C_{\text{CFL}}\,\min\!\left(\frac{\Delta x}{|u|+\sqrt{gh}},\ \frac{\Delta y}{|v|+\sqrt{gh}}\right)
\end{equation}
Por eso el ``paso inicial'' solo es el arranque.

\paragraph{DG2 frente a FV1.} FV1 representa la solución en cada celda con un valor constante (1.er orden). DG2 la representa con un plano (valor medio y dos pendientes, 2.º orden), lo que captura mejor los frentes y el detalle del terreno, pero puede dar profundidades negativas en laderas muy empinadas.

\subsection{Profundidades negativas: regla por volumen}
Las profundidades negativas se ponen en cero y se mide el volumen añadido frente al agua del dominio en cada cuadro:
\begin{equation}
f=\frac{V_{\text{añadido}}}{V_{\text{agua}}}\qquad
\begin{cases}
f<1\% & \text{resultado normal}\\
1\%\le f\le 10\% & \text{aviso: profundidades poco confiables; la dirección del flujo vale}\\
f>10\% & \text{falla; el servidor prueba el otro método}
\end{cases}
\end{equation}

\subsection{Infiltración: número de curva SCS}
Retención potencial $S$ (mm) a partir del número de curva tabulado $CN$ (de WorldCover + grupo de suelo HYSOGs):
\begin{equation}
S_{0.20}=\frac{25\,400}{CN}-254,\qquad S_{0.05}=0.8187\,S_{0.20}^{\,1.15}
\end{equation}
Abstracción inicial $I_a=\lambda S$ con $\lambda=0.05$ (Hawkins et al., 2002), y lluvia neta (escorrentía) acumulada $Q$ para una lluvia acumulada $P$:
\begin{equation}
Q=\begin{cases}\dfrac{(P-\lambda S)^2}{P+(1-\lambda)S} & P>\lambda S\\[6pt] 0 & P\le \lambda S\end{cases}
\end{equation}
Por eso los primeros minutos de lluvia no producen agua en el mapa: se los lleva la abstracción inicial.

\paragraph{Humedad previa (AMC).} Se convierte el $CN$ de condición media (II) a seca (I) o húmeda (III) (Chow et al., 1988):
\begin{equation}
CN_{\text{I}}=\frac{4.2\,CN_{\text{II}}}{10-0.058\,CN_{\text{II}}},\qquad CN_{\text{III}}=\frac{23\,CN_{\text{II}}}{10+0.13\,CN_{\text{II}}}
\end{equation}
En modo automático, la saturación relativa $\theta$ de NISAR SME2 (o SMAP L4) decide la condición: I si $\theta<0.30$, III si $\theta>0.60$, II en otro caso.

\paragraph{Coeficiente de escorrentía (alternativa simple).} Con duración $D$ en horas ($D\ge0.5$):
\begin{equation}
C=\min\!\left(0.92,\ 0.65+0.27\left(1-e^{-D/24}\right)\right),\qquad R_{\text{neta}}=C\cdot R
\end{equation}

\subsection{Lluvia}
\paragraph{Tormenta de diseño (PXR-4).} En cada celda de 0.25° la intensidad máxima anual sigue una distribución GEV cuya ubicación y escala dependen de la duración $D$ (h):
\begin{equation}
\mu=a\,D^{\alpha},\qquad \sigma=b\,D^{\beta},\qquad
i_T=\mu+\sigma\,\frac{1-\left[-\ln\!\left(1-\tfrac{1}{T}\right)\right]^{\kappa}}{\kappa},\quad \kappa=-0.114
\end{equation}
La lluvia total es $P_T=i_T\cdot D$ (mm). Ejemplo en Tuxtla: 100 años y 3 h dan unos 60 mm.

\paragraph{Hietograma de bloques alternos.} Se calcula la lluvia acumulada $P(k\Delta t)$ para bloques de $\Delta t=10$ min, se obtienen los incrementos $\Delta P_k=P(k\Delta t)-P((k-1)\Delta t)$, se ordenan de mayor a menor y se colocan con el mayor al centro y los demás alternando a derecha e izquierda. La suma es exactamente $P(n\Delta t)$.

\paragraph{Lluvia sintética.} Con intensidad media $\bar r$ (total entre duración) y $N$ intervalos (par), se sortean $N/2$ desviaciones $\delta_j\sim U(-b,b)$ con $b=\min(10\ \text{mm/h},\ \bar r)$, se agregan sus opuestas $-\delta_j$ y se barajan con la semilla:
\begin{equation}
r_k=\bar r+\delta_k,\qquad \sum_{k=1}^{N}\delta_k=0\ \Rightarrow\ \text{el total se conserva.}
\end{equation}

\subsection{Alturas y nivel del mar}
\paragraph{Un solo datum vertical (EGM2008).} Con las ondulaciones del geoide $N$ de la NGA:
\begin{align}
H_{\text{EGM2008}}&=h_{\text{WGS84}}-N_{\text{EGM2008}} && \text{(Copernicus GLO-30)}\\
H_{\text{EGM2008}}&=H_{\text{EGM96}}+N_{\text{EGM96}}-N_{\text{EGM2008}} && \text{(NASADEM)}
\end{align}

\paragraph{Barómetro inverso.} Con $\rho=1000$ kg/m³ y $g=9.81$ m/s²:
\begin{equation}
\Delta\eta=\frac{\max(0,\ 1013.25-p)\times100}{\rho g}\ \approx\ 1\ \text{cm por hPa}
\end{equation}

\subsection{Balance de agua}
\begin{equation}
V_{\text{bruta}}=V_{\text{infiltrada}}+V_{\text{neta}},\qquad
V_{\text{neta}}+V_{\text{inicial}}=V_{\text{final}}+V_{\text{salida}}-V_{\text{entrada}},\qquad
V_{\text{final}}=V_{\text{láminas}}+V_{\text{encharcada}}
\end{equation}
Láminas: profundidad $<0.03$ m; encharcada: $\ge0.03$ m. Ejemplo real (Sumidero, 100 años y 3 h): bruta 9078 m³, infiltrada 6285 m³ (69\%), neta 2793 m³, final 2792 m³ (2741 encharcada, 31 en láminas), salida 0.4 m³.

\section{Límites y preguntas probables}
\begin{itemize}
  \item \textbf{¿Está validado?} No. Se buscó un evento observable con lluvia IMERG y agua vista por satélite en la ROI y en la costa, y no hubo uno utilizable entre junio y noviembre de 2026. Las herramientas de validación (precisión, recall, F1, IoU) ya existen.
  \item \textbf{¿Por qué el agua no llena el cañón del río?} El modelo no tiene un río que entre desde fuera del dominio: solo simula la lluvia que cae dentro. Además, el terreno de 30 m no incluye el agua del río.
  \item \textbf{¿Por qué al final sigue mojado?} El modelo no tiene evaporación; las películas delgadas quedan en terreno plano.
  \item \textbf{¿Por qué cambia de DG2 a FV1?} En laderas muy empinadas DG2 genera profundidades negativas; si superan el 10\% del volumen, el servidor usa FV1 y lo avisa.
  \item \textbf{¿Funciona en cualquier lugar?} Sí: todas las fuentes son globales y abiertas. El límite práctico es el tamaño (3000 celdas, 120 s por intento de solver).
  \item \textbf{Ciudad:} el terreno de 30 m no distingue edificios uno por uno, y no hay drenaje pluvial; los resultados urbanos llevan aviso.
\end{itemize}

\section*{Referencias principales}
{\small
\begin{itemize}
  \item Shaw, J., Kesserwani, G., Neal, J., Bates, P., Sharifian, M. K. LISFLOOD-FP 8.0: the new discontinuous Galerkin shallow-water solver for multi-core CPUs and GPUs. \emph{Geosci. Model Dev.} 14, 3577–3602 (2021), doi:10.5194/gmd-14-3577-2021.
  \item Courty, L. G., Wilby, R. L., Hillier, J. K., Slater, L. J. Intensity-duration-frequency curves at the global scale. \emph{Environ. Res. Lett.} 14, 084045 (2019). PXR-4 v2.1.0, doi:10.5281/zenodo.3351812.
  \item Hawkins, R. H., Jiang, R., Woodward, D. E., Hjelmfelt, A. T., Van Mullem, J. A. Runoff curve number method: examination of the initial abstraction ratio. \emph{2nd Federal Interagency Hydrologic Modeling Conference} (2002).
  \item Chow, V. T., Maidment, D. R., Mays, L. W. \emph{Applied Hydrology}. McGraw-Hill (1988).
  \item ESA WorldCover 2021 (CC BY 4.0); HYSOGs250m (ORNL DAAC); Copernicus GLO-30; NASADEM; EOT20; GEBCO\_2026; HydroLAKES v1.0; NGA EGM2008.
\end{itemize}
}

\end{document}
